% GENERATED FILE. Edit canonical lessons, then run npm run generate:books.
% canonical-source-hash: fnv1a64:01f71d276eb4bdd3
% canonical-lessons: JA-C113-hanko, JA-C113-dango, JA-C113-dashi, JA-C113-katsuo, JA-C113-wakame

\chapter{Personal seal, Dumpling, Soup stock, Bonito, Wakame seaweed}
\label{ch:ja-personal-seal-dumpling-soup-stock-bonito-wakame-seaweed}
\hlchaptermodality{113}

\begin{chapteropening}
\textbf{By the end of this chapter:} \emph{I can say a personal seal, a dumpling, soup stock, a bonito and wakame seaweed.}

Complete the last lesson of chapter 113: I can say a personal seal, a dumpling, soup stock, a bonito and wakame seaweed.
\end{chapteropening}

\section[\texorpdfstring{\ja{はんこ} (hanko)}{はんこ (hanko)}]{\ja{はんこ} (hanko) — a personal seal}

\label{lesson:JA-C113-hanko}



\begin{quote}
Before the new one: say the Japanese for work, then the Japanese for a holiday.
\end{quote}

\subsection*{You'll want to know: \ja{はんこ}}
\textbf{\ja{はんこ}} — \emph{hanko} — \textquotedblleft{}a personal seal\textquotedblright{}.

\begin{grammarlens}[title={What you've built}]
One more word for this chapter.
\end{grammarlens}

\subsection*{Your turn}
\begin{itemize}
\raggedright
  \item \emph{Say it:} \emph{hanko}
  \item \emph{Say it:} \emph{hanko}, once more
  \item \emph{Say it:} say \emph{hanko}
  \item \emph{Recall:} say the Japanese for work, then the Japanese for a holiday, then say \emph{hanko} again
\end{itemize}

\begin{tcolorbox}[breakable,colback=teal!4,colframe=teal!35!black,title={Before you move on}]
What does \ja{はんこ} mean? (\textquotedblleft{}A personal seal\textquotedblright{}.) Say it once more.
\end{tcolorbox}
\section[\texorpdfstring{\ja{だんご} (dango)}{だんご (dango)}]{\ja{だんご} (dango) — a dumpling}

\label{lesson:JA-C113-dango}



\begin{quote}
Before the new one: say the Japanese for a holiday, then the Japanese for a personal seal.
\end{quote}

\subsection*{You'll want to know: \ja{だんご}}
\textbf{\ja{だんご}} — \emph{dango} — \textquotedblleft{}a dumpling\textquotedblright{}.

\begin{grammarlens}[title={What you've built}]
One more word for this chapter.
\end{grammarlens}

\subsection*{Your turn}
\begin{itemize}
\raggedright
  \item \emph{Say it:} \emph{dango}
  \item \emph{Say it:} \emph{dango}, once more
  \item \emph{Say it:} say \emph{dango}
  \item \emph{Recall:} say the Japanese for a holiday, then the Japanese for a personal seal, then say \emph{dango} again
\end{itemize}

\begin{tcolorbox}[breakable,colback=teal!4,colframe=teal!35!black,title={Before you move on}]
What does \ja{だんご} mean? (\textquotedblleft{}A dumpling\textquotedblright{}.) Say it once more.
\end{tcolorbox}
\section[\texorpdfstring{\ja{だし} (dashi)}{だし (dashi)}]{\ja{だし} (dashi) — soup stock}

\label{lesson:JA-C113-dashi}



\begin{quote}
Before the new one: say the Japanese for a personal seal, then the Japanese for a dumpling.
\end{quote}

\subsection*{You'll want to know: \ja{だし}}
\textbf{\ja{だし}} — \emph{dashi} — \textquotedblleft{}soup stock\textquotedblright{}.

\begin{grammarlens}[title={What you've built}]
One more word for this chapter.
\end{grammarlens}

\subsection*{Your turn}
\begin{itemize}
\raggedright
  \item \emph{Say it:} \emph{dashi}
  \item \emph{Say it:} \emph{dashi}, once more
  \item \emph{Say it:} say \emph{dashi}
  \item \emph{Recall:} say the Japanese for a personal seal, then the Japanese for a dumpling, then say \emph{dashi} again
\end{itemize}

\begin{tcolorbox}[breakable,colback=teal!4,colframe=teal!35!black,title={Before you move on}]
What does \ja{だし} mean? (\textquotedblleft{}Soup stock\textquotedblright{}.) Say it once more.
\end{tcolorbox}
\section[\texorpdfstring{\ja{かつお} (katsuo)}{かつお (katsuo)}]{\ja{かつお} (katsuo) — a bonito}

\label{lesson:JA-C113-katsuo}



\begin{quote}
Before the new one: say the Japanese for a dumpling, then the Japanese for soup stock.
\end{quote}

\subsection*{You'll want to know: \ja{かつお}}
\textbf{\ja{かつお}} — \emph{katsuo} — \textquotedblleft{}a bonito\textquotedblright{}.

\begin{grammarlens}[title={What you've built}]
One more word for this chapter.
\end{grammarlens}

\subsection*{Your turn}
\begin{itemize}
\raggedright
  \item \emph{Say it:} \emph{katsuo}
  \item \emph{Say it:} \emph{katsuo}, once more
  \item \emph{Say it:} say \emph{katsuo}
  \item \emph{Recall:} say the Japanese for a dumpling, then the Japanese for soup stock, then say \emph{katsuo} again
\end{itemize}

\begin{tcolorbox}[breakable,colback=teal!4,colframe=teal!35!black,title={Before you move on}]
What does \ja{かつお} mean? (\textquotedblleft{}A bonito\textquotedblright{}.) Say it once more.
\end{tcolorbox}
\section[\texorpdfstring{\ja{わかめ} (wakame)}{わかめ (wakame)}]{\ja{わかめ} (wakame) — wakame seaweed}

\label{lesson:JA-C113-wakame}



\begin{quote}
Before the new one: say the Japanese for soup stock, then the Japanese for a bonito.
\end{quote}

\subsection*{You'll want to know: \ja{わかめ}}
\textbf{\ja{わかめ}} — \emph{wakame} — \textquotedblleft{}wakame seaweed\textquotedblright{}.

\begin{grammarlens}[title={What you've built}]
One more word for this chapter.
\end{grammarlens}

\subsection*{Your turn}
\begin{itemize}
\raggedright
  \item \emph{Say it:} \emph{wakame}
  \item \emph{Say it:} \emph{wakame}, once more
  \item \emph{Say it:} say \emph{wakame}
  \item \emph{Recall:} say the Japanese for soup stock, then the Japanese for a bonito, then say \emph{wakame} again
\end{itemize}

\begin{tcolorbox}[breakable,colback=teal!4,colframe=teal!35!black,title={Before you move on}]
What does \ja{わかめ} mean? (\textquotedblleft{}Wakame seaweed\textquotedblright{}.) Say it once more.
\end{tcolorbox}
